\documentclass{article}
\usepackage[T1]{fontenc}
\usepackage[utf8]{inputenc}
\usepackage[polish]{babel}
\usepackage{amsmath}
\usepackage{url}
\usepackage{amssymb}
%{Informatyka stosowana 2020, I st., semestr VI}


\author{
	{1}\\
	{2}\\
	{3}
}

\title{Komputerowe systemy rozpoznawania 2024/2025\\Projekt 1. Klasyfikacja dokumentów tekstowych}
\begin{document}
\maketitle

Opis projektu ma formę artykułu naukowego lub raportu z zadania
badawczego (doświadczalnego/obliczeniowego0 (wg indywidualnych potrzeb związanych np. z
pracą inżynierską/naukową/zawodową). Kolejne sekcje muszą być numerowane i
zatytułowane. Wzory są numerowane, tablice są numerowane i podpisane nad
tablicą, rysunki są numerowane i podpisane pod rysunkiem. Podpis rysunku i
tabeli musi być wyczerpujący (nie ogólnikowy), aby czytelnik nie musiał sięgać do tekstu, aby go
zrozumieć.\\
\indent {\bf Limit stron sprawozdania: 20. Nie zmieniać parametrów tekstu
(czcionki, marginesów, interlinii, itp.). Istotne dla sprawozdania rysunki,
tabele, informacje, kody itp. można umieścić w załącznikach, które nie wliczają
się do limitu stron.}\\
\indent {\bf Wybrane sekcje (rozdziały sprawozdania) są uzupełniane wg wymagań w
opisie Projektu 1. i Harmonogramie Zajęć na WIKAMP KSR jako efekty zadań w~poszczególnych tygodniach}.

\section{Cel projektu}
Zwięzły (2-3 zdania) opis
problemu badawczego/obliczeniowego,  uwzględniający część badawczą i
implementacyjną. {\bf Nie przepisuj literatury ani teorii -- napisz krótko jak
rozumiesz to co masz wykonać: jakie działania, na jakim zbiorze danych (link lub
przypis), jaki jest spodziewany efekt}.\\
\indent Zamieszczony opis (własny, nie skopiowany) zawiera
przypisy do literatury (bibliografii) zamieszczonej na końcu raportu/sprawozdania
zgodnie z~Polską Normą cytowania bibliografii (zob. materiały BG PŁ pt. ,,Bibliografia
załącznikowa'').\\
\noindent {\bf Sekcja uzupełniona jako efekt zadania Tydzień 02 wg Harmonogramu Zajęć
na WIKAMP KSR.}

Celem niniejszego projektu jest implementacja, analiza oraz optymalizacja algorytmu $k$-najbliższych sąsiadów ($k$-NN, z ustaloną wartością początkową $k=5$) przeznaczonego do kategoryzacji dokumentów tekstowych na podstawie wyekstrahowanych cech o charakterze liczbowym i tekstowym.
Klasyfikacja prowadzona jest na korpusie tekstowym \textit{Reuters-21578} \cite{reuters21578}, a zadanie polega na przyporządkowaniu artykułów do jednej z 6 klas reprezentujących kraje (USA, UK, Francja, Kanada, Japonia, Niemcy Zachodnie).
Proces oceny modelu uwzględnia zjawisko braku zbalansowania danych poprzez stratyfikowany podział zbioru na część treningową oraz testową (w proporcji 80/20) i obejmuje zbadanie skuteczności za pomocą metryk Accuracy oraz Macro-F1.


\section{Klasyfikacja nadzorowana metodą $k$-NN.  Ekstrakcja cech, wektory cech}
Krótki opis metody $k$-NN: zasada działania, wymagane parametry wejściowe, format
i~znaczenie wyników/rezultatów. Opis własny z przypisami do literatury -- minimum
teorii potrzebnej do zadania, tak by inżynier innej specjalności zrozumiał dalszy
opis \cite{tadeusiewicz90}. {\bf Nie przepisuj literatury ani teorii -- napisz krótko jak
rozumiesz to co masz wykonać w tym konkretnym przypadku}.\\

Wyekstrahowane cechy liczbowe i tekstowe dokumentów, min. 10 cech, w tym min. 2
o wartościach tekstowych, wszystkie
opisane słownie oraz wzorami, z objaśnieniem oznaczeń i przykładami użycia, do
tego precyzyjny opis możliwych wartości, które przyjmuje dana cecha (ułatwiający
czytelnikowi zrozumienie znaczenia w zadaniu klasyfikacji). Pamiętaj, że wybrane
cechy muszą reprezentować obiekt niezależnie od innych tekstów w tym samym lub w
innym zbiorze. Podaj postać wektora wartości cech po procesie ekstrakcji. Użyte oznaczenia są jednolite w całym
raporcie/sprawozdaniu. \\
\indent Nadmiarowe rysunki i wzory mogą być podane w załącznikach. \\
\noindent {\bf Sekcja uzupełniona jako efekt zadania Tydzień 02 wg Harmonogramu Zajęć
na WIKAMP KSR.}

\subsection{Opis metody $k$-NN w zadaniu kategoryzacji dokumentów}
Algorytm $k$-najbliższych sąsiadów ($k$-NN) to nadzorowana metoda klasyfikacji, polegająca na przypisaniu nowego obiektu do klasy, która najczęściej występuje wśród jego $k$ najbliższych obiektów ze zbioru uczącego (treningowego) \cite{tadeusiewicz90}. W kontekście tego zadania, każdy dokument jest przekształcany w wielowymiarowy wektor wyekstrahowanych cech.
Odległość pomiędzy wektorami (dokumentami) obliczana jest na podstawie zdefiniowanej metryki (np. Euklidesowej czy Czebyszewa dla cech liczbowych i miar podobieństwa napisów dla cech tekstowych). Algorytm przyjmuje jako parametry wejściowe zbiór tekstów, zestaw wag dla każdej z cech oraz liczbę $k$ ($k=5$), a na wyjściu zwraca najbardziej prawdopodobną klasę docelową (kraj, do którego przypisany jest tekst).

\subsection{Ekstrakcja i opis cech wektora}
Dokumenty reprezentowane są za pomocą zbioru zdefiniowanych cech tekstowych i liczbowych, wyekstrahowanych na podstawie zliczeń znaków i słów. Niech $D$ oznacza zbiór znaków występujących w artykule, $L$ oznacza długość dokumentu w znakach (ze spacjami), $W$ to zbiór słów (wyekstrahowanych za pomocą wyrażeń regularnych), $N = |W|$ to liczba słów. Wszystkie zaimplementowane poniżej cechy pozwalają na utworzenie wektora $V$. Poniżej szczegółowo objaśniono stosowane charakterystyki:

\textbf{Cechy liczbowe:}

1. \textbf{CharCountSpec ($c_1$):} Ilość wszystkich znaków w dokumencie bez spacji.
   \begin{equation}
   c_1 = |\{x \in D \mid x \neq \text{' '}\}|
   \end{equation}
   Wartość: $\mathbb{N} \cup \{0\}$.

2. \textbf{WordCountSpec ($c_2$):} Ilość wszystkich słów w danym tekście. Słowo jest zdefiniowane jako zbiór liter bez cyfr.
   \begin{equation}
   c_2 = N
   \end{equation}
   Wartość: $\mathbb{N} \cup \{0\}$.

3. \textbf{AverageWordLengthSpec ($c_3$):} Średnia długość słowa.
   \begin{equation}
   c_3 = \frac{\sum_{w \in W} |w|}{N}
   \end{equation}
   Wartość: $\mathbb{R}_{\ge 0}$. (Gdy brak słów, cecha równa 0).

4. \textbf{LexicalDiversitySpec ($c_4$):} Stosunek ilości unikalnych słów do wszystkich słów w dokumencie. Im większy wynik, tym większy zasób słownictwa w danym tekście. Unikalne słowa badane są bez względu na wielkość znaków. Niech $U$ będzie zbiorem słów z $W$ zredukowanych do małych liter.
   \begin{equation}
   c_4 = \frac{|U|}{N}
   \end{equation}
   Wartość: $[0, 1]$.

5. \textbf{LongWordToOtherWordsRatioSpec ($c_5$):} Stosunek liczby słów bardzo długich (co najmniej 5 liter) do liczby pozostałych słów. Niech $L_w = |\{w \in W \mid |w| \ge 5\}|$.
   \begin{equation}
   c_5 = \begin{cases}
    \frac{L_w}{N - L_w} & \text{dla } N - L_w > 0 \\
    L_w & \text{w p. p.}
   \end{cases}
   \end{equation}
   Wartość: $\mathbb{R}_{\ge 0}$.

6. \textbf{ProperNounRatioSpec ($c_6$):} Współczynnik nazewnictwa własnego – stosunek wyrazów rozpoczynających się wielką literą (z pominięciem tych po znaku końca zdania) do całkowitej ilości słów.
Niech $P$ będzie liczbą takich nazw.
   \begin{equation}
   c_6 = \frac{P}{N}
   \end{equation}
   Wartość: $[0, 1]$.

7. \textbf{UpperToAllCharsRatioSpec ($c_7$):} Stosunek wielkich liter do całkowitej długości tekstu.
   Niech $U_c = |\{x \in D \mid x \text{ jest wielką literą}\}|$.
   \begin{equation}
   c_7 = \frac{U_c}{L}
   \end{equation}
   Wartość: $[0, 1]$.

8. \textbf{UpperToLowerRatioSpec ($c_8$):} Stosunek ilości wielkich liter do małych liter.
   Niech $M_c = |\{x \in D \mid x \text{ jest małą literą}\}|$.
   \begin{equation}
   c_8 = \begin{cases}
    \frac{U_c}{M_c} & \text{dla } M_c > 0 \\
    U_c & \text{w p. p.}
   \end{cases}
   \end{equation}
   Wartość: $\mathbb{R}_{\ge 0}$.

9. \textbf{AlnumToOtherCharsRatioSpec ($c_9$):} Stosunek liczby znaków alfanumerycznych (litery i cyfry) do liczby pozostałych znaków (interpunkcja, znaki białe, znaki specjalne).
   Niech $A_c = |\{x \in D \mid x \text{ jest literą lub cyfrą}\}|$, a $O_c = L - A_c$.
   \begin{equation}
   c_9 = \begin{cases}
    \frac{A_c}{O_c} & \text{dla } O_c > 0 \\
    A_c & \text{w p. p.}
   \end{cases}
   \end{equation}
   Wartość: $\mathbb{R}_{\ge 0}$.

\textbf{Cechy tekstowe:}

10. \textbf{ProperNounSpec ($t_1$):} Zwraca pierwszą napotkaną w tekście nazwę własną (wyraz zaczynający się wielką literą). Jeśli takiej nie ma, przyjmuje wartość tekstową "None".
    \begin{equation}
    t_1 = w_i \in W, \text{gdzie } w_i \text{ zaczyna się od wielkiej litery oraz } \forall_{j < i} w_j \text{ jest małą literą}
    \end{equation}
    Wartość: Ciąg znaków (String).

11. \textbf{RarestRepeatedWordSpec ($t_2$):} Wyraz o najmniejszej ilości powtórzeń, jednak pojawiający się przynajmniej dwukrotnie w artykule (redukcja do małych liter). W przypadku remisów wybierane jest słowo wcześniejsze leksykograficznie. Brak takich wyrazów skutkuje przyjęciem wartości "None".
    Niech $C(w)$ to liczba wystąpień słowa $w$ w tekście.
    \begin{equation}
    t_2 = \underset{w \in U, C(w) \ge 2}{\mathrm{argmin}}\ C(w)
    \end{equation}
    Wartość: Ciąg znaków (String).

\textbf{Postać wektora cech:}

Na podstawie zaprezentowanych specyfikacji każdy tekst zostaje przekształcony na wektor $V$ postaci:
\begin{equation}
V = [c_1, c_2, c_3, c_4, c_5, c_6, c_7, c_8, c_9, t_1, t_2]
\end{equation}
Przykład wektora dla zdania ,,ALFA alfa 123 !!! kot kot Dom dom DOM'':
\begin{equation}
V_{przyklad} = [30.0, 8.0, 3.25, 0.5, 0.0, 0.0, 0.25, 0.6, 1.4, \text{''ALFA''}, \text{''alfa''}]
\end{equation}

\section{Miary jakości klasyfikacji}
Miary jakości klasyfikacji (Accuracy, Precision,
Recall, F1). We wprowadzeniu zaprezentować minimum teorii potrzebnej do realizacji
zadania, tak by inżynier innej specjalności zrozumiał dalszy opis. Należy podać {\bf konkretne wzory miar użyte w tym eksperymencie oraz krótko
opisać ich znaczenie i zakresy przyjmowanych wartości. Należy podać przykładowe
wartości każdej miary. Nie przepisuj
teorii, ale podaj link/przypis i opisz jak rozumiesz jej zastosowanie w tym konkretnym
zadaniu}. \\
\indent Stosowane wzory, oznaczenia z objaśnieniami znaczenia symboli użytych w
doświadczeniu. Oznaczenia jednolite w obrębie całego sprawozdania.  Opis zawiera przypisy do bibliografii zgodnie z
Polską Normą, (zob. materiały BG PŁ).\\
\noindent {\bf Sekcja uzupełniona jako efekt zadania Tydzień 03 wg Harmonogramu Zajęć
na WIKAMP KSR.}


\section{Metryki i miary podobieństwa tekstów w klasyfikacji}
Wzory, znaczenia i opisy symboli zastosowanych metryk z
przykładami. Wzory, opisy i znaczenia miar
podobieństwa tekstów zastosowanych w obliczaniu metryk dla wektorów cech z
przykładami dla każdej miary \cite{niewiadomski08}.  Oznaczenia jednolite w obrębie całego sprawozdania.  {\bf Podaj metryki i miary
podobieństwa nie z literatury (te wystarczy zacytować linkiem), ale konkretne ich
postaci stosowane w zadaniu. Jakie zakresy wartości przyjmują te miary i
metryki, co oznaczają ich wartości? Podaj przykładowe wartości dla przykładowych wektorów cech}. \\
\noindent {\bf Sekcja uzupełniona jako efekt zadania Tydzień 04 wg Harmonogramu Zajęć
na WIKAMP KSR.}

\section{Wyniki klasyfikacji dla różnych parametrów wejściowych}
Wstępne wyniki miary Accuracy dla próbnych klasyfikacji na ograniczonym zbiorze tekstów (podać parametry i kryteria
wyboru wg punktów 3.-8. z opisu Projektu 1.).
\noindent {\bf Sekcja uzupełniona jako efekt zadania Tydzień 05 wg Harmonogramu Zajęć
na WIKAMP KSR.}


\section{Dyskusja, wnioski, sprawozdanie końcowe}
We wnioskach z badania skomentować wyłącznie wpływ parametrów i cech na jakość procesu klasyfikacji. Innymi słowy – wnioski to to, co wynika z przeprowadzonych obliczeń, a nie to co wiadomo z literatury.
Znaczenie kolejnych eksperymentów wg punktów 2.-8. opisu projektu 1.  Każdorazowo
podane parametry, dla których przeprowadzana eksperyment.
Wykresy (np. słupkowe) i tabele wyników
obowiązkowe, dokładnie opisane w ,,captions'' (tytułach), konieczny opis osi i
jednostek wykresów oraz kolumn i wierszy tabel.\\

{**Ewentualne wyniki realizacji punktu 9. opisu Projektu 1., czyli ,,na ocenę 5.0'' i ich porównanie do wyników z
części obowiązkowej**.Dokładne interpretacje uzyskanych wyników w zależności od parametrów klasyfikacji
opisanych w punktach 3.-8 opisu Projektu 1.
Omówić i wyjaśnić napotkane problemy (jeśli były). Każdy wniosek/problem powinien mieć poparcie
w przeprowadzonych eksperymentach (odwołania do konkretnych wyników: wykresów,
tabel). \\
\underline{Dla końcowej oceny jest to najważniejsza sekcja} sprawozdania, gdyż prezentuje poziom
zrozumienia rozwiązywanego problemu.\\

** Możliwości kontynuacji prac w obszarze systemów rozpoznawania, zwłaszcza w kontekście pracy inżynierskiej,
magisterskiej, naukowej, itp. **\\

\noindent {\bf Sekcja uzupełniona jako efekt zadań Tydzień 05 i Tydzień 06 wg Harmonogramu Zajęć
na WIKAMP KSR.}


\section{Braki w realizacji projektu 1.}
Wymienić wg opisu Projektu 1. wszystkie niezrealizowane obowiązkowe elementy projektu, ewentualnie
podać merytoryczne (ale nie czasowe) przyczyny tych braków.


\begin{thebibliography}{0}
\bibitem{tadeusiewicz90} R. Tadeusiewicz: Rozpoznawanie obrazów, PWN, Warszawa, 1991.
\bibitem{niewiadomski08} A. Niewiadomski, Methods for the Linguistic Summarization of Data: Applications of Fuzzy Sets and Their Extensions, Akademicka Oficyna Wydawnicza EXIT, Warszawa, 2008.
\bibitem{reuters21578} David D. Lewis, \textit{Reuters-21578 Text Categorization Collection Distribution 1.0}, UCI Machine Learning Repository, 1997. Dostępne online: \url{http://kdd.ics.uci.edu/databases/reuters21578/reuters21578.html}
\end{thebibliography}

Literatura zawiera wyłącznie źródła recenzowane i/lub o potwierdzonej wiarygodności,
możliwe do weryfikacji i cytowane w sprawozdaniu.
\end{document}
